\documentclass[a4paper,12pt]{article}

\usepackage[spanish]{babel}
\usepackage[utf8]{inputenc}
\usepackage{graphicx}
\usepackage{hyperref}
\usepackage{listings}
\usepackage{xcolor}
\usepackage{geometry}

\geometry{margin=2.5cm}

\title{Simulador de Bolsa Web con Análisis de Portafolio}
\date{\today}

\lstset{
language=JavaScript,
basicstyle=\ttfamily\small,
keywordstyle=\color{blue},
stringstyle=\color{red},
commentstyle=\color{green!60!black}
}

\begin{document}

% -------- HOJA DE PRESENTACION --------

\begin{titlepage}

\centering

\vspace*{2cm}

{\Large \textbf{Segundo Proyecto}}\\[1cm]

{\LARGE \textbf{Simulador de Bolsa Web con Análisis de Portafolio}}\\[3cm]

\begin{tabular}{rl}

\textbf{Nombre del alumno:} & Williams Espinosa López \\

\textbf{Matrícula:} & 251185 \\

\textbf{Cuatrimestre y grupo:} & 4° “B” \\

\textbf{Docente:} & Sirgei García Ballinas \\

\textbf{Asignatura:} & Cálculo Integral \\

\end{tabular}

\vfill

{\large \today}

\end{titlepage}

% -------- DOCUMENTO --------

\section{Introducción}

Los mercados financieros son sistemas complejos donde los inversionistas compran y venden activos como acciones, bonos y derivados. Para comprender su comportamiento es útil utilizar herramientas de simulación.

Este proyecto consiste en el desarrollo de un simulador de bolsa basado en tecnologías web que permite:

\begin{itemize}
\item Simular la compra y venta de acciones
\item Analizar el rendimiento de un portafolio
\item Aplicar estrategias de trading
\item Evaluar el riesgo de inversión
\item Utilizar predicciones mediante inteligencia artificial
\item Almacenar operaciones en una base de datos
\end{itemize}

El sistema se implementa utilizando HTML, CSS y JavaScript, integrando además servicios externos para obtener datos financieros reales.

\section{Arquitectura del Sistema}

La aplicación está compuesta por tres capas principales:

\begin{itemize}
\item \textbf{Frontend}: Interfaz web desarrollada en HTML, CSS y JavaScript.
\item \textbf{API de datos financieros}: Alpha Vantage para obtener precios reales de acciones.
\item \textbf{Base de datos}: Supabase con PostgreSQL para almacenar operaciones.
\end{itemize}

\section{Simulación de Portafolios}

Un portafolio es un conjunto de activos financieros que posee un inversor. El simulador permite crear un portafolio virtual donde el usuario puede comprar y vender acciones.

El valor del portafolio se calcula mediante la fórmula:

\[
Valor\ del\ Portafolio = \sum_{i=1}^{n} Cantidad_i \times Precio_i
\]

También se calcula el retorno de la inversión (ROI):

\[
ROI = \frac{Valor\ actual - Inversión\ inicial}{Inversión\ inicial} \times 100
\]

El simulador permite visualizar estos resultados en tiempo real mediante gráficos.

\section{Simulación de Estrategias de Trading}

El sistema implementa distintas estrategias de trading para analizar posibles escenarios de inversión.

\subsection{Buy and Hold}

La estrategia Buy and Hold consiste en comprar un activo y mantenerlo durante un largo periodo esperando que su valor aumente.

\subsection{Momentum}

La estrategia Momentum se basa en la idea de que los activos que han mostrado una tendencia alcista continuarán subiendo en el corto plazo.

\section{Análisis de Riesgo}

El riesgo de una inversión se mide normalmente mediante la volatilidad.

La volatilidad se calcula utilizando la desviación estándar:

\[
\sigma = \sqrt{\frac{\sum (x_i - \mu)^2}{n}}
\]

donde:

\begin{itemize}
\item $x_i$ representa cada precio
\item $\mu$ es la media de los precios
\item $n$ es el número de observaciones
\end{itemize}

Una mayor volatilidad indica mayor riesgo en la inversión.

\section{Uso de Aprendizaje Automático}

El simulador incluye un módulo de predicción basado en aprendizaje automático. Este módulo genera estimaciones del precio futuro de un activo utilizando modelos probabilísticos simples.

En implementaciones más avanzadas podrían utilizarse algoritmos como:

\begin{itemize}
\item Redes neuronales
\item Regresión lineal
\item Modelos ARIMA
\end{itemize}

\section{Integración con Bases de Datos}

Para almacenar las operaciones de compra y venta se utiliza Supabase, una plataforma basada en PostgreSQL.

Cada operación incluye:

\begin{itemize}
\item Tipo de operación
\item Símbolo de la acción
\item Cantidad
\item Precio
\item Fecha
\end{itemize}

Esto permite mantener un historial de operaciones y analizar el rendimiento del portafolio.

\section{Obtención de Datos Financieros}

Para obtener datos reales del mercado se utiliza la API Alpha Vantage.

Ejemplo de consulta:

\begin{lstlisting}
https://www.alphavantage.co/query
?function=GLOBAL_QUOTE
&symbol=AAPL
&apikey=TU_API_KEY
\end{lstlisting}

Esta consulta devuelve el precio actual de una acción.

\section{Tecnologías Utilizadas}

\begin{itemize}
\item HTML5
\item CSS3
\item JavaScript
\item Chart.js
\item Alpha Vantage API
\item Supabase (PostgreSQL)
\end{itemize}

\section{Conclusión}

El simulador desarrollado demuestra cómo las tecnologías web modernas pueden utilizarse para crear herramientas de análisis financiero.

La integración de APIs externas, bases de datos y algoritmos de simulación permite construir aplicaciones educativas que ayudan a comprender el funcionamiento de los mercados financieros.

\end{document}